\section{The full higher even symplectic normal-core family}
\label{sec:higher-even-symplectic}

\noindent\textbf{Status.} Complete new authored proof. Root and independent peer
/root/g34\_eightfolds have each whole-reviewed the proof and reported
PASS; root also read the full author certificate and reran it successfully.
The mathematical version reviewed by both has SHA-256
\texttt{72e295707440f2e06e39658d2b3528f61d9587227e99e6f18196b075515dbc3c}. The elementary construction in D-HES-02
replaces the unverified Gluck--Magaard (1998) input in the archived
conditional draft. That unverified source claim is not used here:
the only character-ratio input is the explicitly verified published
unipotent bound (HES.3). The finite tables and their precise provenance
are retained in D-HES-F.

\subsection*{D-HES-01. Elementary degree bounds and verified inputs}

Put $S_n=\operatorname{Sp}_{2n}(2)$ for $n\ge4$. These groups are simple,
centerless, have trivial Schur multiplier and have no outer automorphisms.
The exact all-$n\ge4$ multiplier and automorphism statements occur in
\cite[author p.~25]{HES:GM2006}; the exceptional-multiplier list in
\cite[Chapter~5, Table~1, numbered p.~35]{BI:HHM} is consistent with this.
These classification facts do not require the fourth-moment hypothesis
used elsewhere in the former paper.

Every nontrivial irreducible complex representation of $S_4$ has degree
at least $35$; for $S_5$ the minimum is $155$. These are complete finite
table assertions already certified in D-HES-F. Only its two ATLAS
endpoints $S_4,S_5$ are essential here: the independently checked
K\"ohler $S_6$ table is redundant for this all-rank theorem.
For $n\ge5$, every
nontrivial irreducible representation of $S_n$ consequently has degree
\[
                     a\ge\max\{155,2^{n-1}-1\}.
\tag{HES.1}
\]
Indeed it is faithful because $S_n$ is simple. Restriction to the block
subgroup $S_5$ has a nontrivial constituent, proving the first bound.
For the second embed $\operatorname{GL}_n(2)$ in $S_n$ by
$A\mapsto\operatorname{diag}(A,A^{-\mathsf T})$. Its subgroup
\[
 U_0=\left\{\begin{pmatrix}1&v\\0&I_{n-1}\end{pmatrix}:
                                  v\in\mathbf F_2^{n-1}\right\}
\]
is elementary abelian, and its normalizer contains
$\operatorname{GL}_{n-1}(2)$ transitively permuting all its nontrivial
characters. Faithfulness makes one nontrivial weight occur, hence all
$2^{n-1}-1$ occur with equal positive multiplicity. This proves (HES.1).

We use two further known inputs with their precise scopes. Gow's theorem
states that every element of a symplectic group in characteristic two is
a product of two involutions, hence is conjugate to its inverse
\cite[theorem on p.~583; proof, Section~3, pp.~589--591]{HES:Gow1981}.
All ordinary irreducible characters of $S_n$ are therefore real valued.
Gluck's unipotent character-ratio estimate, stated with its definitions in
the published survey \cite[Theorem~2.4(i), author pp.~2--4]{HES:Liebeck2017},
gives for our groups
\[
       |\chi(u)|\le\tfrac34\chi(1)
          \qquad(1\ne u\text{ unipotent},\quad 1\ne\chi\in\operatorname{Irr}(S_n)).
\tag{HES.3}
\]
The nonunipotent part of the same source has a different bound and is not
silently strengthened. The original reference for (HES.3) is Gluck (1995),
as specified in that survey.

\subsection*{D-HES-02. An elementary parabolic argument}

For a finite-order operator $x$ on a complex space of dimension $a$, write
\[
           \nu(x)=a-\max_\lambda\dim\ker(x-\lambda I).
\]

\begin{lemma}[D-HES-02a: fixed-point-free elementary subgroup]
Let $h\in S_n$ have odd prime order $p$, and suppose that $h$ stabilizes
a nonzero totally isotropic subspace of the natural symplectic space
$W=\mathbf F_2^{2n}$. Then there is a nonzero $h$-normalized elementary
abelian $2$-subgroup $U\le S_n$ with $C_U(h)=1$ and $|U|\ge4$.
\end{lemma}

\begin{proof}
Let $R$ be such an isotropic subspace. First suppose that $h|_R$ is
nontrivial, and choose a basis $r_1,\ldots,r_s$ of $R$. Write $b$ for the
symplectic pairing. For a symmetric matrix $B=(b_{ij})$ over $\mathbf F_2$
define
\[
             N_B(x)=\sum_{i,j}b_{ij}b(r_j,x)r_i.
\]
Since $R$ is isotropic, $N_BN_C=0$ for all $B,C$. Symmetry of $B$ and
characteristic two imply
$b(N_Bx,y)+b(x,N_By)=0$, while $b(N_Bx,N_By)=0$.
Thus $1+N_B$ is symplectic, and these operators form an elementary
abelian subgroup, canonically the additive group of symmetric matrices.
If $A$ is the matrix of $h|_R$, conjugation by $h$ sends $B$ to
$ABA^{\mathsf T}$. This action is nontrivial: for a rank-one matrix
$B=vv^{\mathsf T}$, the image is $(Av)(Av)^{\mathsf T}$, and over
$\mathbf F_2$ the map $v\mapsto vv^{\mathsf T}$ is injective (its diagonal
recovers $v$). Nonidentity $A$ therefore moves one of these matrices.

Since $p$ is odd, this $\mathbf F_2\langle h\rangle$-module is semisimple.
Its sum of nontrivial simple summands is a nonzero invariant subspace
without fixed vectors. The corresponding subgroup is the desired $U$.

Now suppose $h$ fixes $R$ pointwise. Choose $0\ne e\in R$, and choose
$f_0$ with $b(e,f_0)=1$. Put
$f=\sum_{j=0}^{p-1}h^jf_0$. Then $h(f)=f$ and $b(e,f)=p=1$ in
$\mathbf F_2$. Consequently
$W=\langle e,f\rangle\perp W_0$ is an $h$-stable orthogonal decomposition,
and $h|_{W_0}$ is nontrivial, since otherwise $h=1$.
For $v\in W_0$ and $c\in\mathbf F_2$ define the symplectic operator
\[
 q(v,c)e=e,\qquad
 q(v,c)w=w+b(w,v)e\ (w\in W_0),\qquad
 q(v,c)f=f+v+ce.
\]
Direct multiplication gives
\[
       q(v,c)q(w,d)=q(v+w,c+d+b(v,w)).
\tag{HES.4}
\]
In characteristic two $b(v,w)=b(w,v)$ and $b(v,v)=0$.
Thus this group $Q$ is abelian of exponent two, of order $2^{2n-1}$.
Moreover $h q(v,c)h^{-1}=q(hv,c)$, so the action on $Q$ is nontrivial.
Applying the same semisimple-module argument to the elementary abelian
group $Q$ gives a nonzero invariant subgroup $U$ with $C_U(h)=1$.

Finally a one-dimensional vector space over $\mathbf F_2$ has trivial
automorphism group. A nonzero fixed-point-free module for an odd-order
element must have dimension at least two; hence $|U|\ge4$ in both cases.
No invariant isotropic complement, parabolic character formula, or
unverified estimate is used in this construction.
\end{proof}

\begin{lemma}[D-HES-02b: all-element support]
For $n\ge5$, every nontrivial irreducible complex representation $\rho$
of $S_n$ of degree $a$ satisfies
\[
                   \nu(\rho(x))\ge a/8\qquad(1\ne x\in S_n).
\tag{HES.5}
\]
For $n=4$, the complete finite table gives the stronger bound $a/5$.
\end{lemma}

\begin{proof}
We first treat an odd-prime-order $h$ in a proper parabolic subgroup;
equivalently it stabilizes a nonzero totally isotropic subspace. Use $U$
from D-HES-02a, and put $t=\dim V^U$. Every nonidentity element of $U$
is an involution, hence unipotent. Averaging (HES.3) gives
\[
       t=\frac1{|U|}\sum_{u\in U}\chi(u)
        \le a\left(\frac34+\frac1{4|U|}\right)
        \le\frac{13}{16}a.
\tag{HES.6}
\]
The dual action on the character group $U^*$ is also fixed-point-free
away from its identity, because $h-1$ is invertible on $U$ and on its
dual. Every nontrivial character orbit has length $p$. On the sum of
weight spaces belonging to such an orbit, $h$ cyclically permutes equal
dimensional spaces, with $h^p=1$. Each $p$-th root therefore occurs with
exactly one $p$-th of that orbit-sum dimension. On $V^U$ we need only the
trivial upper bound $t$. It follows that every $h$-eigenspace has
dimension at most $t+(a-t)/p$, so
\[
       \nu(\rho(h))\ge(1-1/p)(a-t)
                   \ge\frac23\frac3{16}a=\frac18a.
\tag{HES.7}
\]

Next suppose an odd-prime-order $h$ belongs to no proper parabolic.
Decompose its natural space $W$ into primary summands. A summand for
an irreducible polynomial not equal to its reciprocal is totally
isotropic: it pairs only with the distinct reciprocal summand. Such a
summand is forbidden. A nonzero $1$-eigenspace is nondegenerate symplectic
and contains an invariant isotropic line, so it too is forbidden.
Each remaining irreducible polynomial is self-reciprocal of degree $2d$.
On its isotypic summand set $E=\mathbf F_{2^{2d}}$ and
$E_0=\mathbf F_{2^d}$. Adjointness of the symplectic pairing induces the
involution $z\mapsto\bar z=z^{2^d}$ on $E$. The trace pairing defines a
nondegenerate Hermitian form $H$ by
\[
              b(av,w)=\operatorname{Tr}_{E/\mathbf F_2}(aH(v,w))
                         \quad(a\in E).
\]
The adjoint identity and symmetry of $b$ give
$H(av,bw)=a\bar b H(v,w)$ and $H(w,v)=\overline{H(v,w)}$.
If the $E$-dimension were at least two, diagonalization and surjectivity
of the finite-field norm $E^\times\to E_0^\times$ would give a nonzero
isotropic vector. Its entire $E$-line would be totally isotropic for $b$,
contradicting the absence of an invariant isotropic subspace.
Every such primary polynomial therefore has multiplicity one.

The centralizer is consequently a product of norm-one scalar groups of
orders $2^{d_i}+1$, with $\sum_i d_i=n$. Hence
\[
                   |C_{S_n}(h)|=\prod_i(2^{d_i}+1)\le3^n.
\tag{HES.8}
\]
Character column orthogonality and (HES.1) imply
\[
        \frac{|\chi(h)|}{a}
        \le\frac{3^{n/2}}{\max\{155,2^{n-1}-1\}}<\frac34.
\tag{HES.9}
\]
For $n=5,6$ compare directly with $155$. For $n=7$ the squared inequality
is $16\cdot3^7=34992<35721=9\cdot63^2$.
For subsequent $n$, write $x=2^{n-1}\ge64$ and observe
$(2x-1)^2-3(x-1)^2=x^2+2x-2>0$. This proves the required induction.
All nonidentity powers of $h$ generate the same cyclic group and are
again elliptic, so (HES.9) holds for every one of them. The Fourier
formula for any $p$-th-root eigenspace bounds its dimension by
\[
             \frac1p\left(a+\sum_{j=1}^{p-1}|\chi(h^j)|\right)
                 \le a\left(\frac34+\frac1{4p}\right)\le\frac56a.
\]
Thus the elliptic case has the stronger support bound $a/6$.

If $h$ has order two, (HES.3) gives directly
$\nu(\rho(h))=(a-|\chi(h)|)/2\ge a/8$.
For arbitrary $1\ne x$, choose a nonidentity prime-order power $h=x^k$.
Every $x$-eigenspace lies in an $h$-eigenspace, so its projective support
is no smaller. This proves (HES.5). Finally D-HES-F gives
$|\chi(x)|\le3a/5$ for $S_4$; if an eigenvalue has multiplicity $m$, then
$|\chi(x)|\ge2m-a$, yielding $\nu\ge a/5$.
\end{proof}

\subsection*{D-HES-03. Full normalizers and the actual period image}

\begin{theorem}[D-HES-03: unconditional all-rank exclusion]
In the signature-$(3,3)$ sixfold-source problem, an irreducible inducing
representation whose faithful image has an irreducible normal subgroup
which is a product of groups $\operatorname{Sp}_{2n_i}(2)$, $n_i\ge4$,
cannot have an actual period image of dimension nine. Arbitrary
normalizing scalars, compatible permutations of factors, and inducing
character fields larger than the original imaginary quadratic field are
allowed. No determinant or reflection assumption is needed.
\end{theorem}

\begin{proof}
Let $J\le\operatorname{GL}(V)$ be the inducing image, $e=\dim V$, and
$E$ the specified irreducible normal subgroup. Its centerless simple
factors make it an actual direct product, and
$V=\bigotimes_iV_i$ with every factor nontrivial. A normalizer operator
permutes isomorphic simple factors. After removing this permutation,
its automorphisms on each factor are inner, since their outer groups
are trivial. Schur's lemma makes it a scalar times a tensor of core
operators. The trivial Schur multipliers also rule out an omitted
projective-cover ambiguity.

For a factor-preserving tensor, eigenvalue multiplicity distributions
convolve on roots of unity, and the maximum normalized multiplicity
cannot exceed that of a nonscalar factor. A nontrivial cycle permuting
$l$ equal-dimensional factors can be conjugated so all twists lie on
one coordinate, and their product can be diagonalized. In the resulting
weighted permutation of $a^l$ basis words there are $a$ constant words
and all remaining orbits have length at least two. A specified
eigenvalue thus has multiplicity at most
$a+(a^l-a)/2\le3a^l/4$. Additional tensor factors cannot increase that
normalized bound. D-HES-02 therefore gives projective support at least
$e/8$ on every nonscalar $J$-operator. Nonidentity scalars have actual
moving rank $e$. In particular every effective nonidentity inertia has
actual moving rank at least $e/8$.

Gow's theorem makes all the factor modules self-dual. Thus $V$ has a
nondegenerate $E$-invariant bilinear form with a unique invariant line,
and $J$ acts on that line by a multiplier $\mu:J\to\mathbf C^\times$.
This multiplier is nontrivial. Otherwise every $J$-operator would have
reciprocal eigenvalue pairs and hence real trace, contradicting inclusion
of the original imaginary quadratic field in the inducing character
field. The field inclusion follows from the rational-induction argument
of D-BI-01 and does not assume equality of the fields.

The actual intermediate-cover factorization of D-BI-01 identifies the
six-dimensional multiplicity variation with the parabolic cohomology
of this inducing representation. For intermediate genus $h$ and
nonidentity effective inertias $g_1,\ldots,g_r$ it gives
\[
                6=2(h-1)e+\sum_{j=1}^r\operatorname{rank}(g_j-I).
\tag{HES.10}
\]
Here $e\ge35$, so $h\ge2$ is impossible. If $h=1$ the sum is six,
whence $r\le6$ and pointed-curve moduli dimension is at most six
(one when $r=0$). The actual finite-cover factorization then bounds
period-image dimension below nine. Thus a nine-dimensional image would
require $h=0$, $r\ge12$ and moving-rank sum $2e+6$.

The generating inertia tuple has product one, so at least two entries
have nontrivial multiplier. Their $1$-eigenspaces are totally isotropic
for a nondegenerate bilinear form and each has dimension at most $e/2$.
Those entries cost at least $e$ in total; the remaining at least ten
cost at least $e/8$ each. Consequently
\[
                  2e+6\ge e+10e/8=9e/4,
                         \qquad e\le24,
\tag{HES.11}
\]
contradicting $e\ge35$. This uses an actual variation factorization and
a necessary branch budget, not a dimension inequality as a sufficiency
test for dominance.
\end{proof}

\paragraph{Verification boundary.}
The finite subgroup constructions in D-HES-02 are explicit linear
algebra, proved above. A separate exact program checks representative
matrices, group laws, the weight-orbit averaging arithmetic, and the
symbolic all-rank inequalities; it is supplementary to these proofs.
No claim to have retrieved or verified Gluck--Magaard (1998), Theorem~3.1,
is made. The former conditional argument is retained as an archive,
and its missing source is no longer a premise of this theorem.
